% Complete reviewed Unit 9; editorial master: chapters/ch09/answer-key.md
% Markdown SHA-256: 28e21ea9437f1d6289487267c5247ae14700bf4aeceda6c183c608d531dd5d1d
% One-use initial assembly followed by complete content and page review.
\newcommand{\runningtitle}{Unit 9. Answer key}
% Layout inherited from the reviewed previous course; set a running title per unit.
% This is a small hand-maintained TeX layout, not a conversion pipeline.
\RequirePackage{fix-cm}
\documentclass[11pt,twoside]{article}
\usepackage{fontspec}
\setmainfont{Cambria}[
  Renderer=HarfBuzz,
  BoldFont=cambriab.ttf,
  ItalicFont=cambriai.ttf,
  BoldItalicFont=cambriaz.ttf
]
\setmonofont{lmmono10-regular.otf}[BoldFont=lmmonolt10-bold.otf]
\usepackage[paperwidth=165mm,paperheight=241mm,inner=21mm,outer=19mm,
  top=22mm,bottom=20mm,headheight=14pt,headsep=7mm,footskip=10mm]{geometry}
\usepackage{xcolor,array,booktabs,longtable,tabularx}
\usepackage{enumitem,fancyhdr}
\usepackage[unicode,hidelinks]{hyperref}
\definecolor{ink}{HTML}{000000}
\definecolor{muted}{HTML}{4C5558}
\setlength{\parindent}{0pt}
\setlength{\parskip}{3.5pt}
\AtBeginDocument{\fontsize{11}{14.3}\selectfont}
\setlength{\tabcolsep}{4pt}
\renewcommand{\arraystretch}{1.10}
\setlist[itemize]{leftmargin=1.3em,label=\textbullet,itemsep=2pt,topsep=3pt,parsep=0pt}
\setlist[enumerate]{leftmargin=1.8em,itemsep=3pt,topsep=4pt,parsep=0pt}
\widowpenalty=10000
\clubpenalty=10000
\emergencystretch=1em
\raggedbottom
\pagestyle{fancy}
\fancyhf{}
\providecommand{\runningtitle}{Old English}
\fancyhead[LE,RO]{\small\itshape \runningtitle}
\fancyhead[LO,RE]{\small Old English}
\fancyfoot[LE,RO]{\small\thepage}
\renewcommand{\headrulewidth}{0.3pt}
\newcommand{\booktitle}[2]{
  {\fontsize{9}{11}\selectfont\color{muted}#1\par}
  \vspace{3pt}
  {\fontsize{20}{24}\selectfont\color{ink}#2\par}
  \vspace{6pt}
}
\newcommand{\source}[1]{\par{\footnotesize\color{muted}#1\par}}
\newcommand{\readingtext}[1]{
  \par\addvspace{5pt}
  \begingroup\leftskip=1.5em
    \fontsize{12}{16}\selectfont #1\par
  \endgroup\addvspace{5pt}
}
\makeatletter
\renewcommand{\section}{\@startsection{section}{1}{\z@}
  {-2.4ex \@plus -0.6ex \@minus -0.2ex}{0.8ex \@plus 0.2ex}
  {\normalfont\fontsize{12}{15}\selectfont\bfseries\color{ink}}}
\renewcommand{\subsection}{\@startsection{subsection}{2}{\z@}
  {-1.8ex \@plus -0.5ex \@minus -0.2ex}{0.5ex \@plus 0.2ex}
  {\normalfont\fontsize{11}{14}\selectfont\bfseries\color{ink}}}
\makeatother

\hyphenpenalty=10000
\exhyphenpenalty=10000
\emergencystretch=3em
\newcommand{\itempara}[2]{\par\addvspace{3pt}\noindent\hangindent=1.8em\hangafter=1\makebox[1.8em][l]{#1.}#2\par}
\newcommand{\reserve}[1]{\par\penalty-100\begingroup\dimen0=\pagegoal\advance\dimen0 by-\pagetotal\ifdim\dimen0<#1\newpage\fi\endgroup}
\begin{document}
\addfontfeatures{Ligatures=TeXOff}
\raggedright
\booktitle{UNIT 9}{Answer key}

% mdparagraph:P001:begin
These models use the course's selected forms. Other documented forms, word orders and translations can be valid. Preserve the participants, number, time and force of the actual task, and distinguish a grammatical answer from one that meets an explicitly requested construction. The Old English answers and explanations are new teaching material; Cook's English prompts retain their original wording. [1]
% mdparagraph:P001:end

% mdparagraph:P002:begin
For a requested report or dependent question, use the stated mood. An indicative content clause can assert a fact; a subjunctive report does not declare its content false. A past reporting verb does not impose an automatic tense change on every subordinate clause. The infinitive after commanding or perceiving has no person ending: its actor and its own object have separate roles even when both are accusative.
% mdparagraph:P002:end

% mdparagraph:P003:begin
In paradigm tasks, give every labeled cell, including identical forms. The ivory, living-creatures and 365-days expressions stay supplied, as do the source reading's local helps; they are not tests of unaided recall.
% mdparagraph:P003:end

\reserve{7\baselineskip}
% mdparagraph:P004:begin
\section*{F. Complete forms and their use}
% mdparagraph:P004:end

\par\begin{minipage}{\linewidth}\raggedright
% item: NEW-U09-F01
% key: KEY-NEW-U09-F01
% mdparagraph:P005:begin
\textbf{F01. Sunu.}
% mdparagraph:P005:end

% mdtable:T01:begin
\par\begingroup
\fontsize{10.5}{13.2}\selectfont
\setlength{\LTleft}{0pt}\setlength{\LTright}{\fill}
\noindent
\begin{tabular}{@{}>{\raggedright\arraybackslash}p{39.46mm}>{\raggedright\arraybackslash}p{39.46mm}>{\raggedright\arraybackslash}p{39.46mm}@{}}
\toprule
\textbf{Case} & \textbf{Singular} & \textbf{Plural} \\
\midrule
Nominative & sunu & suna \\
Genitive & suna & suna \\
Dative & suna & sunum \\
Accusative & sunu & suna \\
\bottomrule
\end{tabular}
\par\endgroup

% mdtable:T01:end
% mdparagraph:P006:begin
\textbf{Target and interpretation:} The complete selected u-stem paradigm has dative singular \textbf{suna}, contrasting with nominative/\allowbreak{}accusative \textbf{sunu}. Plural \textbf{suna} serves three cases, but dative plural is \textbf{sunum}. Cook documents \textbf{-u/\allowbreak{}-o} variants outside genitives and dative plural, and later a-stem \textbf{sunes/\allowbreak{}sunas}; these remain valid where applicable, distinct from this requested table. All vowels in the selected forms are short. [2]
% mdparagraph:P006:end

\end{minipage}\par\addvspace{4pt}

\par\begin{minipage}{\linewidth}\raggedright
% item: NEW-U09-F02
% key: KEY-NEW-U09-F02
% mdparagraph:P007:begin
\textbf{F02. Bycgan: complete indicative.}
% mdparagraph:P007:end

% mdtable:T02:begin
\par\begingroup
\fontsize{10.5}{13.2}\selectfont
\setlength{\LTleft}{0pt}\setlength{\LTright}{\fill}
\noindent
\begin{tabular}{@{}>{\raggedright\arraybackslash}p{44.51mm}>{\raggedright\arraybackslash}p{36.93mm}>{\raggedright\arraybackslash}p{36.93mm}@{}}
\toprule
\textbf{Person and number} & \textbf{Present} & \textbf{Preterite} \\
\midrule
First singular & bycge & bohte \\
Second singular & bygest & bohtest \\
Third singular & bygeð & bohte \\
First plural & bycgað & bohton \\
Second plural & bycgað & bohton \\
Third plural & bycgað & bohton \\
\bottomrule
\end{tabular}
\par\endgroup

% mdtable:T02:end
% mdparagraph:P008:begin
\textbf{Target and interpretation:} Give all twelve cells, with irregular dental preterite \textbf{bohte}. Present \textbf{bycge/\allowbreak{}bycgað} contrasts with selected \textbf{bygest/\allowbreak{}bygeð}; do not preserve \textbf{cg} mechanically in every present form. Preterite first and third singular coincide, while second singular adds \textbf{-st}. Shared plural forms still represent three labeled persons. A regular past invented from \textbf{bycgan} misses the lexical stem.
% mdparagraph:P008:end

\end{minipage}\par\addvspace{4pt}

\par\begin{minipage}{\linewidth}\raggedright
% item: NEW-U09-F03
% key: KEY-NEW-U09-F03
% mdparagraph:P009:begin
\textbf{F03. Commands.}
% mdparagraph:P009:end

% mdtable:T03:begin
\par\begingroup
\fontsize{10.5}{13.2}\selectfont
\setlength{\LTleft}{0pt}\setlength{\LTright}{\fill}
\noindent
\begin{tabular}{@{}>{\raggedright\arraybackslash}p{26.9mm}>{\raggedright\arraybackslash}p{49.64mm}>{\raggedright\arraybackslash}p{41.84mm}@{}}
\toprule
\textbf{Verb} & \textbf{Second singular: one addressee} & \textbf{Second plural: two or more} \\
\midrule
rōwan & rōw & rōwað \\
bycgan & byge & bycgað \\
\bottomrule
\end{tabular}
\par\endgroup

% mdtable:T03:end
% mdparagraph:P010:begin
\textbf{Target and interpretation:} Give both commands of each verb. \textbf{Rōw} lacks indicative second-singular \textbf{-st}; \textbf{byge} is the selected singular command, not first-singular indicative \textbf{bycge}. The plural commands use the present plural stems. No separate dual command is required. These models omit the subject; the separately attested postposed-pronoun construction is outside this target.
% mdparagraph:P010:end

\end{minipage}\par\addvspace{4pt}

\par\begin{minipage}{\linewidth}\raggedright
% item: NEW-U09-F04
% key: KEY-NEW-U09-F04
% mdparagraph:P011:begin
\textbf{F04. Rīce and geðeahtend.}
% mdparagraph:P011:end

% mdtable:T04:begin
\par\begingroup
\fontsize{10.5}{13.2}\selectfont
\setlength{\LTleft}{0pt}\setlength{\LTright}{\fill}
\noindent
\begin{tabular}{@{}>{\raggedright\arraybackslash}p{38.56mm}>{\raggedright\arraybackslash}p{41.26mm}>{\raggedright\arraybackslash}p{38.56mm}@{}}
\toprule
\textbf{Case} & \textbf{Rīce: singular} & \textbf{Rīce: plural} \\
\midrule
Nominative & rīce & rīcu \\
Genitive & rīces & rīca \\
Dative & rīce & rīcum \\
Accusative & rīce & rīcu \\
\bottomrule
\end{tabular}
\par\endgroup

% mdtable:T04:end
% mdtable:T05:begin
\par\begingroup
\fontsize{10.5}{13.2}\selectfont
\setlength{\LTleft}{0pt}\setlength{\LTright}{\fill}
\noindent
\begin{tabular}{@{}>{\raggedright\arraybackslash}p{29.37mm}>{\raggedright\arraybackslash}p{45.82mm}>{\raggedright\arraybackslash}p{43.18mm}@{}}
\toprule
\textbf{Case} & \textbf{Geðeahtend: singular} & \textbf{Geðeahtend: plural} \\
\midrule
Nominative & geðeahtend & geðeahtend \\
Genitive & geðeahtendes & geðeahtendra \\
Dative & geðeahtende & geðeahtendum \\
Accusative & geðeahtend & geðeahtend \\
\bottomrule
\end{tabular}
\par\endgroup

% mdtable:T05:end
% mdparagraph:P012:begin
\textbf{Target and interpretation:} Give two complete eight-cell systems. In \textbf{rīce}, lexical \textbf{-e} gives way to \textbf{rīcu/\allowbreak{}rīca/\allowbreak{}rīcum}; Cook also permits inserted \textbf{i/\allowbreak{}e} in these plurals, including \textbf{rīciu/\allowbreak{}rīcia}. \textbf{Geðeahtend} has selected endingless nominative/\allowbreak{}accusative plural, genitive plural \textbf{-ra}, dative plural \textbf{-um}. Documented nominative/\allowbreak{}accusative plural \textbf{-e} and rarer \textbf{-as} are alternatives. This agent noun does not require an inflecting participle's adjective endings. [2]
% mdparagraph:P012:end

\end{minipage}\par\addvspace{4pt}

\par\begin{minipage}{\linewidth}\raggedright
% item: NEW-U09-F05
% key: KEY-NEW-U09-F05
% mdparagraph:P013:begin
\textbf{F05. Kingdoms and counsellors.}
% mdparagraph:P013:end

% mdparagraph:P014:begin
\textbf{Ðāra geðeahtendra rīcu.} — The kingdoms of the counsellors.
% mdparagraph:P014:end

% mdparagraph:P015:begin
\textbf{Ic sęlle ðone fisc ðǣm geðeahtendum.} — I give the fish to the counsellors.
% mdparagraph:P015:end

% mdparagraph:P016:begin
\textbf{Target and interpretation:} \textbf{Geðeahtendra} is genitive plural dependent on nominative plural \textbf{rīcu}, with agreeing \textbf{ðāra}. The head has no separate article, as requested. \textbf{Geðeahtendum} is dative plural recipient of \textbf{sęlle}, with dative \textbf{ðǣm}; accusative singular \textbf{ðone fisc} is the thing given. Substituting genitive for the recipient would lose that relation. Documented \textbf{rīciu} is a valid plural variant.
% mdparagraph:P016:end

\end{minipage}\par\addvspace{4pt}

\reserve{7\baselineskip}
% mdparagraph:P017:begin
\section*{A. Content and reported statements}
% mdparagraph:P017:end

\par\begin{minipage}{\linewidth}\raggedright
% item: NEW-U09-A01
% key: KEY-NEW-U09-A01
% mdparagraph:P018:begin
\textbf{A01. Three functions of ðæt.}
% mdparagraph:P018:end

\begin{itemize}
% mdparagraph:P019:begin
\item \textbf{(a) Ðæt scip is gōd:} \textbf{ðæt} is the article/\allowbreak{}demonstrative, nominative neuter singular, agreeing with subject \textbf{scip}.
% mdparagraph:P019:end
% mdparagraph:P020:begin
\item \textbf{(b) Ic wāt, ðæt ðæt scip gōd is:} the first \textbf{ðæt} is an invariant content conjunction, without case. The second is the nominative neuter singular article/\allowbreak{}demonstrative agreeing with the content clause's subject \textbf{scip}.
% mdparagraph:P020:end
% mdparagraph:P021:begin
\item \textbf{(c) Ðæt scip, ðæt ðū hæfst, is gōd:} the first \textbf{ðæt} again agrees with nominative subject \textbf{scip}. The relative \textbf{ðæt} is accusative neuter singular, the object of \textbf{hæfst} inside its own clause; \textbf{ðū} is that clause's subject.
% mdparagraph:P021:end
\end{itemize}

% mdparagraph:P022:begin
\textbf{Target and interpretation:} Distinguish agreement with a noun from a relative pronoun's own clause role. The relative and its antecedent agree in gender and number here, but their cases differ. Identical spelling neither makes all five occurrences conjunctions nor gives the conjunction an accusative case.
% mdparagraph:P022:end

\end{minipage}\par\addvspace{4pt}

\par\begin{minipage}{\linewidth}\raggedright
% item: NEW-U09-A02
% key: KEY-NEW-U09-A02
% mdparagraph:P023:begin
\textbf{A02. An asserted fact.}
% mdparagraph:P023:end

% mdparagraph:P024:begin
\textbf{Ic wāt, ðæt ðā scipu ǣmtig sindon.} — I know that the ships are empty.
% mdparagraph:P024:end

% mdparagraph:P025:begin
\textbf{Target and interpretation:} The ships' emptiness is asserted as fact: indicative plural \textbf{sindon} agrees with \textbf{ðā scipu}. Predicate \textbf{ǣmtig} is selected strong neuter nominative plural, despite its endingless appearance. Singular \textbf{is} would lose the number. A contextual subjunctive can be valid elsewhere but misses the requested indicative; \textbf{ðæt} alone does not demand subjunctive.
% mdparagraph:P025:end

\end{minipage}\par\addvspace{4pt}

\par\begin{minipage}{\linewidth}\raggedright
% item: NEW-U09-A03
% key: KEY-NEW-U09-A03
% mdparagraph:P026:begin
\textbf{A03. A present report.}
% mdparagraph:P026:end

% mdparagraph:P027:begin
\textbf{Se fiscere sægð, ðæt hē twēgen fiscas hæbbe.} — The fisherman says that he has two fish.
% mdparagraph:P027:end

% mdparagraph:P028:begin
\textbf{Target and interpretation:} The fisherman remains the possessor when quoted \textbf{ic} becomes reported \textbf{hē}. \textbf{Twēgen fiscas} is accusative masculine plural and \textbf{hæbbe} the requested present subjunctive. Retaining \textbf{ic} would ordinarily make the reporting speaker the possessor. Indicative \textbf{hæfð} can be defensible reporting but misses this target; subjunctive does not imply a lie.
% mdparagraph:P028:end

\end{minipage}\par\addvspace{4pt}

\par\begin{minipage}{\linewidth}\raggedright
% item: NEW-U09-A04
% key: KEY-NEW-U09-A04
% mdparagraph:P029:begin
\textbf{A04. A simultaneous past report.}
% mdparagraph:P029:end

% mdparagraph:P030:begin
\textbf{Se mann sægde, ðæt hē twēgen fiscas hæfde.} — The man said that he had two fish.
% mdparagraph:P030:end

% mdparagraph:P031:begin
\textbf{Target and interpretation:} The same man possessed the fish at the time of his past speech. Requested \textbf{hæfde} is preterite subjunctive; the same spelling also occurs in the indicative, so the task and construction establish its analysis. The example teaches simultaneous past report, not a rule forcing a tense change after every past reporting verb.
% mdparagraph:P031:end

\end{minipage}\par\addvspace{4pt}

\par\begin{minipage}{\linewidth}\raggedright
% item: NEW-U09-A05
% key: KEY-NEW-U09-A05
% mdparagraph:P032:begin
\textbf{A05. Grammar and the requested mood.}
% mdparagraph:P032:end

% mdparagraph:P033:begin
\textbf{Hē sægð, ðæt hē gōd sīe.} — He says that he is good.
% mdparagraph:P033:end

% mdparagraph:P034:begin
This meets the requested subjunctive target. The given indicative sentence is not necessarily ungrammatical.
% mdparagraph:P034:end

% mdparagraph:P035:begin
\textbf{Target and interpretation:} Give \textbf{sīe} for the requested present subjunctive singular, while recognizing that indicative \textbf{is} can be grammatical. Correct the target mismatch without declaring all indicative reports invalid or saying that subjunctive proves the reported claim false.
% mdparagraph:P035:end

\end{minipage}\par\addvspace{4pt}

\reserve{7\baselineskip}
% mdparagraph:P036:begin
\section*{Q. Dependent questions}
% mdparagraph:P036:end

\par\begin{minipage}{\linewidth}\raggedright
% item: NEW-U09-Q01
% key: KEY-NEW-U09-Q01
% mdparagraph:P037:begin
\textbf{Q01. Subject and recipient inside a question.}
% mdparagraph:P037:end

% mdparagraph:P038:begin
\textbf{(a) Ic wāt, hwā ðone fisc hæbbe.} — I know who has the fish. \textbf{Hwā} is nominative, the subject of \textbf{hæbbe}.
% mdparagraph:P038:end

% mdparagraph:P039:begin
\textbf{(b) Ic wāt, hwǣm ðū ðone fisc sęlle.} — I know to whom you give the fish. \textbf{Hwǣm} is dative recipient; \textbf{ðū} is subject and the fish the accusative object.
% mdparagraph:P039:end

% mdparagraph:P040:begin
\textbf{Target and interpretation:} The question words take their cases within their own clauses: \textbf{hwā} is subject of \textbf{hæbbe}, while \textbf{hwǣm} receives what \textbf{ðū} gives. Matrix \textbf{wāt} does not make both accusative. The selected subjunctives express dependent questions; they do not change “who has” into “whom you have.”
% mdparagraph:P040:end

\end{minipage}\par\addvspace{4pt}

\par\begin{minipage}{\linewidth}\raggedright
% item: NEW-U09-Q02
% key: KEY-NEW-U09-Q02
% mdparagraph:P041:begin
\textbf{Q02. A dependent object question.}
% mdparagraph:P041:end

% mdparagraph:P042:begin
\textbf{Saga mē, hwæt ðū hæbbe.} — Tell me what you have.
% mdparagraph:P042:end

% mdparagraph:P043:begin
\textbf{Target and interpretation:} Singular \textbf{saga} takes dative recipient \textbf{mē}. Accusative \textbf{hwæt} remains the thing possessed; \textbf{ðū} remains the addressee and governs subjunctive \textbf{hæbbe}. Indicative \textbf{hæfst} can be valid in another context but misses the requested mood. No English-like auxiliary is added to the dependent question.
% mdparagraph:P043:end

\end{minipage}\par\addvspace{4pt}

\par\begin{minipage}{\linewidth}\raggedright
% item: NEW-U09-Q03
% key: KEY-NEW-U09-Q03
% mdparagraph:P044:begin
\textbf{Q03. A dependent yes/\allowbreak{}no question.}
% mdparagraph:P044:end

% mdparagraph:P045:begin
\textbf{Saga mē, hwæðer se mann twēgen fiscas hæbbe.} — Tell me whether the man has two fish.
% mdparagraph:P045:end

% mdparagraph:P046:begin
\textbf{Target and interpretation:} Invariant \textbf{hwæðer} introduces whether the possession is true; it does not ask which of two men possesses the fish. Subject \textbf{se mann} takes singular subjunctive \textbf{hæbbe}, with accusative \textbf{twēgen fiscas}. Inflecting \textbf{hwæðer} as an adjective changes the function; indicative \textbf{hæfð} misses the explicit mood target.
% mdparagraph:P046:end

\end{minipage}\par\addvspace{4pt}

\par\begin{minipage}{\linewidth}\raggedright
% item: NEW-U09-Q04
% key: KEY-NEW-U09-Q04
% mdparagraph:P047:begin
\textbf{Q04. Which of two, or whether?}
% mdparagraph:P047:end

% mdparagraph:P048:begin
\textbf{(a) Hwæðer mann is ðīn gefēra?} — Which of the two men is your companion? \textbf{Hwæðer} is nominative masculine singular, agreeing with \textbf{mann}.
% mdparagraph:P048:end

% mdparagraph:P049:begin
\textbf{(b) Saga mē, hwæðer se mann ðīn gefēra sīe.} — Tell me whether the man is your companion. Here \textbf{hwæðer} is invariant whether.
% mdparagraph:P049:end

% mdparagraph:P050:begin
\textbf{Target and interpretation:} The first question selects one man from the stated two; the second asks whether an identification holds. “Which man” does not translate the second question. The invariant conjunction has no masculine nominative agreement, although its spelling coincides with that adjective cell.
% mdparagraph:P050:end

\end{minipage}\par\addvspace{4pt}

\par\begin{minipage}{\linewidth}\raggedright
% item: NEW-U09-Q05
% key: KEY-NEW-U09-Q05
% mdparagraph:P051:begin
\textbf{Q05. Whom you help.}
% mdparagraph:P051:end

% mdparagraph:P052:begin
\textbf{Ic wāt, hwǣm ðū helpe.} — I know whom you help.
% mdparagraph:P052:end

% mdparagraph:P053:begin
\textbf{Target and interpretation:} \textbf{Helpan} governs dative \textbf{hwǣm} inside the question. \textbf{Helpe} is present subjunctive second singular, without indicative stem mutation. Documented dative \textbf{hwām} is also valid. Accusative \textbf{hwone} misses the government; indicative \textbf{hilpst} misses the requested mood.
% mdparagraph:P053:end

\end{minipage}\par\addvspace{4pt}

\reserve{7\baselineskip}
% mdparagraph:P054:begin
\section*{I. Infinitive complements and object predicates}
% mdparagraph:P054:end

\par\begin{minipage}{\linewidth}\raggedright
% item: NEW-U09-I01
% key: KEY-NEW-U09-I01
% mdparagraph:P055:begin
\textbf{I01. Who carries the fish?}
% mdparagraph:P055:end

% mdparagraph:P056:begin
\textbf{(a) Ic wille ðone fisc beran.} — I want, or am willing, to carry the fish.
% mdparagraph:P056:end

% mdparagraph:P057:begin
\textbf{(b) Ic hāte hine ðone fisc beran.} — I command him to carry the fish.
% mdparagraph:P057:end

% mdparagraph:P058:begin
\textbf{Target and interpretation:} The finite verbs agree with \textbf{ic}; the carrier is the speaker in (a), the expressed accusative \textbf{hine} in (b). The fish is what is carried. An infinitive does not gain a person ending from its actor. These are separately selected lexical frames, not a rule for all verbs.
% mdparagraph:P058:end

\end{minipage}\par\addvspace{4pt}

\par\begin{minipage}{\linewidth}\raggedright
% item: NEW-U09-I02
% key: KEY-NEW-U09-I02
% mdparagraph:P059:begin
\textbf{I02. Reporting the command.}
% mdparagraph:P059:end

% mdparagraph:P060:begin
\textbf{Se fiscere hēt hīe ðone fisc beran.} — The fisherman commanded them to carry the fish.
% mdparagraph:P060:end

% mdparagraph:P061:begin
\textbf{Target and interpretation:} \textbf{Hēt} is past command; accusative plural \textbf{hīe} names the two companions. Their original direct command \textbf{berað} becomes plain infinitive \textbf{beran}. The finite verb agrees with the singular fisherman, not the plural actors. Retaining \textbf{berað} would leave a direct-command form inside the requested reported-command construction.
% mdparagraph:P061:end

\end{minipage}\par\addvspace{4pt}

\par\begin{minipage}{\linewidth}\raggedright
% item: NEW-U09-I03
% key: KEY-NEW-U09-I03
% mdparagraph:P062:begin
\textbf{I03. Perceiving an action.}
% mdparagraph:P062:end

% mdparagraph:P063:begin
\textbf{Ic gesēo ðā męnn ðone fisc beran.} — I see the men carrying the fish.
% mdparagraph:P063:end

% mdparagraph:P064:begin
\textbf{Ðā męnn} is the accusative plural actor of \textbf{beran}; \textbf{ðone fisc} is that infinitive's accusative singular object.
% mdparagraph:P064:end

% mdparagraph:P065:begin
\textbf{Target and interpretation:} Finite first-singular \textbf{gesēo} agrees with \textbf{ic}. The accusative men are the observed actors; the accusative fish is the infinitive's object. Both phrases have the same case but different functions. A participle or separate finite clause would miss the prescribed infinitive construction.
% mdparagraph:P065:end

\end{minipage}\par\addvspace{4pt}

\par\begin{minipage}{\linewidth}\raggedright
% item: NEW-U09-I04
% key: KEY-NEW-U09-I04
% mdparagraph:P066:begin
\textbf{I04. Good as an object predicate.}
% mdparagraph:P066:end

% mdparagraph:P067:begin
\textbf{God geworhte ðone mann gōdne.} — God made the man good.
% mdparagraph:P067:end

% mdparagraph:P068:begin
\textbf{God geworhte ðā męnn gōde.} — God made the men good.
% mdparagraph:P068:end

% mdparagraph:P069:begin
\textbf{Target and interpretation:} Strong \textbf{gōdne/\allowbreak{}gōde} is accusative masculine singular/\allowbreak{}plural, agreeing with the object. The article within that object phrase does not make this separate predicate weak. Nominative \textbf{gōd} or weak \textbf{gōdan} misses the selected agreeing predicate; the finite verb stays singular because God is the subject in both sentences.
% mdparagraph:P069:end

\end{minipage}\par\addvspace{4pt}

\par\begin{minipage}{\linewidth}\raggedright
% item: NEW-U09-I05
% key: KEY-NEW-U09-I05
% mdparagraph:P070:begin
\textbf{I05. Naming and commanding.}
% mdparagraph:P070:end

% mdparagraph:P071:begin
\textbf{Hē hēt ðone mann cyning.} — He called, or named, the man king.
% mdparagraph:P071:end

% mdparagraph:P072:begin
\textbf{Hē hēt ðone mann gān.} — He commanded the man to go.
% mdparagraph:P072:end

% mdparagraph:P073:begin
\textbf{Target and interpretation:} The complements identify the two senses of \textbf{hātan}; identical \textbf{hēt} does not settle the meaning alone. \textbf{Cyning} is a predicate accusative naming the object, whereas \textbf{gān} is a plain infinitive whose actor is the man. Neither complement should be analyzed as the other grammatical category.
% mdparagraph:P073:end

\end{minipage}\par\addvspace{4pt}

\reserve{7\baselineskip}
% mdparagraph:P074:begin
\section*{W. Present unreal wishes}
% mdparagraph:P074:end

\par\begin{minipage}{\linewidth}\raggedright
% item: NEW-U09-W01
% key: KEY-NEW-U09-W01
% mdparagraph:P075:begin
\textbf{W01. A ship wished for now.}
% mdparagraph:P075:end

% mdparagraph:P076:begin
\textbf{Hæfde ic ān scip!} — Would that I had one ship!
% mdparagraph:P076:end

% mdparagraph:P077:begin
\textbf{Target and interpretation:} The ship is wished for now: first-singular preterite subjunctive \textbf{hæfde}, with accusative neuter singular \textbf{ān scip}. \textbf{Ēalā, hæfde ic ān scip!} is also valid. A statement about a formerly possessed ship changes the meaning; a volitional “I want” construction would miss the explicitly requested subjunctive.
% mdparagraph:P077:end

\end{minipage}\par\addvspace{4pt}

\par\begin{minipage}{\linewidth}\raggedright
% item: NEW-U09-W02
% key: KEY-NEW-U09-W02
% mdparagraph:P078:begin
\textbf{W02. One spelling, two settings.}
% mdparagraph:P078:end

% mdparagraph:P079:begin
\textbf{(a) Ic hæfde ān scip.} — I had one ship.
% mdparagraph:P079:end

% mdparagraph:P080:begin
\textbf{(b) Hæfde ic ān scip!} — Would that I had one ship!
% mdparagraph:P080:end

% mdparagraph:P081:begin
\textbf{Target and interpretation:} In (a), \textbf{hæfde} is first-singular preterite indicative; in (b), first-singular preterite subjunctive expressing a present unreal wish. The shared form alone cannot prove mood or force past time. Initial verb position suits this wish but is not a universal test for subjunctive.
% mdparagraph:P081:end

\end{minipage}\par\addvspace{4pt}

\reserve{7\baselineskip}
% mdparagraph:P082:begin
\section*{Cook Group II, Exercise 9}
% mdparagraph:P082:end

% mdparagraph:P083:begin
The four complete Cook sets retain their original prompts. A model may select one defensible interpretation where English leaves number or reference open; the notes identify consequential alternatives. [3]
% mdparagraph:P083:end

\par\begin{minipage}{\linewidth}\raggedright
% item: EX-II-09-01
% key: KEY-EX-II-09-01
% mdparagraph:P084:begin
\textbf{Cook II.9.1.} Sparks fly.
% mdparagraph:P084:end

% mdparagraph:P085:begin
\textbf{Fȳrspearcan flēogað.}
% mdparagraph:P085:end

% mdparagraph:P086:begin
\textbf{Target and interpretation:} Plural \textbf{fȳrspearcan} is the subject of present plural \textbf{flēogað}. The weak noun's \textbf{-an} can have other analyses elsewhere, but is nominative plural here. Singular \textbf{flīehð} would change the sparks into one spark; no article is required for the indefinite plural.
% mdparagraph:P086:end

\end{minipage}\par\addvspace{4pt}

\par\begin{minipage}{\linewidth}\raggedright
% item: EX-II-09-02
% key: KEY-EX-II-09-02
% mdparagraph:P087:begin
\textbf{Cook II.9.2.} The house is empty.
% mdparagraph:P087:end

% mdparagraph:P088:begin
\textbf{Ðæt hūs is ǣmtig.}
% mdparagraph:P088:end

% mdparagraph:P089:begin
\textbf{Target and interpretation:} Nominative neuter singular \textbf{ðæt hūs} takes singular \textbf{is} and strong predicate \textbf{ǣmtig}. The article agrees with the noun; it does not make a separate predicate adjective weak. The sentence states emptiness now, not in the past.
% mdparagraph:P089:end

\end{minipage}\par\addvspace{4pt}

\par\begin{minipage}{\linewidth}\raggedright
% item: EX-II-09-03
% key: KEY-EX-II-09-03
% mdparagraph:P090:begin
\textbf{Cook II.9.3.} The ships are empty.
% mdparagraph:P090:end

% mdparagraph:P091:begin
\textbf{Ðā scipu sindon ǣmtig.}
% mdparagraph:P091:end

% mdparagraph:P092:begin
\textbf{Target and interpretation:} Nominative neuter plural \textbf{ðā scipu} takes plural \textbf{sindon} and selected strong neuter plural \textbf{ǣmtig}. The predicate shares its spelling with the singular in II.9.2; the article, noun and finite verb show the number change. Do not infer singular solely from the adjective's zero ending.
% mdparagraph:P092:end

\end{minipage}\par\addvspace{4pt}

\par\begin{minipage}{\linewidth}\raggedright
% item: EX-II-09-04
% key: KEY-EX-II-09-04
% mdparagraph:P093:begin
\textbf{Cook II.9.4.} The merchant sells needles, awls, fish-hooks, and plowshares.
% mdparagraph:P093:end

% mdparagraph:P094:begin
\textbf{Se mangere cīepð nǣdla, ǣla, anglas, and sulhscearas.}
% mdparagraph:P094:end

% mdparagraph:P095:begin
\textbf{Target and interpretation:} Singular \textbf{cīepð} means sells, with four accusative plural objects: feminine \textbf{nǣdla/\allowbreak{}ǣla}, masculine \textbf{anglas/\allowbreak{}sulhscearas}. \textbf{Angel} loses its medial vowel in \textbf{anglas}. \textbf{Bygeð} means buys and would reverse the merchant's action. The objects' plurality does not make the singular merchant's verb plural.
% mdparagraph:P095:end

\end{minipage}\par\addvspace{4pt}

\par\begin{minipage}{\linewidth}\raggedright
% item: EX-II-09-05
% key: KEY-EX-II-09-05
% mdparagraph:P096:begin
\textbf{Cook II.9.5.} The blacksmith makes plowshares in his smithy.
% mdparagraph:P096:end

% mdparagraph:P097:begin
\textbf{Se smið gewyrcð sulhscearas on his smiððan.}
% mdparagraph:P097:end

% mdparagraph:P098:begin
\textbf{Target and interpretation:} Singular \textbf{gewyrcð} takes accusative plural \textbf{sulhscearas}. Location \textbf{on his smiððan} contains dative singular weak feminine \textbf{smiððe}. Invariant \textbf{his} identifies the male possessor, rather than agreeing with the smithy. The man works in the smithy; this is not motion into it.
% mdparagraph:P098:end

\end{minipage}\par\addvspace{4pt}

\par\begin{minipage}{\linewidth}\raggedright
% item: EX-II-09-06
% key: KEY-EX-II-09-06
% mdparagraph:P099:begin
\textbf{Cook II.9.6.} The oil and the wine are good.
% mdparagraph:P099:end

% mdparagraph:P100:begin
\textbf{Se ęle and ðæt wīn sindon gōd.}
% mdparagraph:P100:end

% mdparagraph:P101:begin
\textbf{Target and interpretation:} The coordinated subjects take plural \textbf{sindon}. Masculine \textbf{ęle} and neuter \textbf{wīn} have their own articles, \textbf{se/\allowbreak{}ðæt}. Selected predicate \textbf{gōd} is neuter plural for this mixed inanimate group. Do not copy masculine plural \textbf{gōde} mechanically from one conjunct; the unit teaches a normal pattern without claiming exceptionless historical agreement.
% mdparagraph:P101:end

\end{minipage}\par\addvspace{4pt}

\par\begin{minipage}{\linewidth}\raggedright
% item: EX-II-09-07
% key: KEY-EX-II-09-07
% mdparagraph:P102:begin
\textbf{Cook II.9.7.} The fisherman rows over the sea.
% mdparagraph:P102:end

% mdparagraph:P103:begin
\textbf{Se fiscere rēwð ofer sǣ.}
% mdparagraph:P103:end

% mdparagraph:P104:begin
\textbf{Target and interpretation:} Singular \textbf{fiscere} takes third-singular present \textbf{rēwð}. \textbf{Ofer sǣ} gives the crossing. The selected extent/\allowbreak{}route analysis can be accusative, but \textbf{sǣ} has identical accusative and dative singular spellings, so its surface alone cannot establish case. \textbf{Rōwað} would not be this singular indicative.
% mdparagraph:P104:end

\end{minipage}\par\addvspace{4pt}

\par\begin{minipage}{\linewidth}\raggedright
% item: EX-II-09-08
% key: KEY-EX-II-09-08
% mdparagraph:P105:begin
\textbf{Cook II.9.8.} The carpenter supports his son.
% mdparagraph:P105:end

% mdparagraph:P106:begin
\textbf{Se trēowwyrhta āfētt his sunu.}
% mdparagraph:P106:end

% mdparagraph:P107:begin
\textbf{Target and interpretation:} The carpenter supports or maintains his son: singular \textbf{āfētt}, invariant possessor \textbf{his}, accusative singular \textbf{sunu}. Genitive/\allowbreak{}dative \textbf{suna} would miss the son's object role. The separately supplied \textbf{fēt} is not substituted by assuming the prefix can simply be removed.
% mdparagraph:P107:end

\end{minipage}\par\addvspace{4pt}

\par\begin{minipage}{\linewidth}\raggedright
% item: EX-II-09-09
% key: KEY-EX-II-09-09
% mdparagraph:P108:begin
\textbf{Cook II.9.9.} The tailor bought the needle.
% mdparagraph:P108:end

% mdparagraph:P109:begin
\textbf{Se sēamere bohte ðā nǣdle.}
% mdparagraph:P109:end

% mdparagraph:P110:begin
\textbf{Target and interpretation:} Past singular \textbf{bohte}, not present \textbf{bygeð}, expresses bought. \textbf{Ðā nǣdle} is accusative feminine singular, the needle bought; \textbf{ðā} is not invariably plural. The irregular preterite must be learned, rather than invented as a regular past from \textbf{bycgan}.
% mdparagraph:P110:end

\end{minipage}\par\addvspace{4pt}

\par\begin{minipage}{\linewidth}\raggedright
% item: EX-II-09-10
% key: KEY-EX-II-09-10
% mdparagraph:P111:begin
\textbf{Cook II.9.10.} The bellows blow.
% mdparagraph:P111:end

% mdparagraph:P112:begin
\textbf{Ðā bęlgas blāwað.}
% mdparagraph:P112:end

% mdparagraph:P113:begin
\textbf{Target and interpretation:} The selected plural \textbf{ðā bęlgas}, referring to the bellows, takes present plural \textbf{blāwað}. A documented single-instrument reading with singular \textbf{se bęlg blǣwð} can be defensible with consistent number and agreement; English final \textbf{-s} alone does not settle Old English number.
% mdparagraph:P113:end

\end{minipage}\par\addvspace{4pt}

\par\begin{minipage}{\linewidth}\raggedright
% item: EX-II-09-11
% key: KEY-EX-II-09-11
% mdparagraph:P114:begin
\textbf{Cook II.9.11.} The hammer smites,
% mdparagraph:P114:end

% mdparagraph:P115:begin
\textbf{Sēo slęcg sliehð.}
% mdparagraph:P115:end

% mdparagraph:P116:begin
\textbf{Target and interpretation:} Feminine singular \textbf{sēo slęcg} takes \textbf{sliehð}, in the explicitly taught absolute sense strike or make a stroke. No object need be invented. \textbf{Bīett}, from \textbf{bēatan}, is a defensible hammer-striking alternative supported by Cook's discussion. The English source's final comma is preserved above; the new model can end with a full stop.
% mdparagraph:P116:end

\end{minipage}\par\addvspace{4pt}

\par\begin{minipage}{\linewidth}\raggedright
% item: EX-II-09-12
% key: KEY-EX-II-09-12
% mdparagraph:P117:begin
\textbf{Cook II.9.12.} Let us seek the kingdom of God.
% mdparagraph:P117:end

% mdparagraph:P118:begin
\textbf{Sēcen wē rīce Godes.}
% mdparagraph:P118:end

% mdparagraph:P119:begin
\textbf{Target and interpretation:} First-plural present subjunctive \textbf{sēcen} with \textbf{wē} expresses the inclusive exhortation. \textbf{Rīce} is accusative singular, the kingdom sought; \textbf{Godes} is genitive. \textbf{Godes rīce} is also a defensible order. Plural command \textbf{sēcað} addresses others and would not itself express the requested “let us.”
% mdparagraph:P119:end

\end{minipage}\par\addvspace{4pt}

\par\begin{minipage}{\linewidth}\raggedright
% item: EX-II-09-13
% key: KEY-EX-II-09-13
% mdparagraph:P120:begin
\textbf{Cook II.9.13.} He will not live with the blacksmith.
% mdparagraph:P120:end

% mdparagraph:P121:begin
\textbf{Hē ne wille libban mid ðǣm smiðe.}
% mdparagraph:P121:end

% mdparagraph:P122:begin
\textbf{Target and interpretation:} \textbf{Ne wille} expresses refusal or unwillingness, followed by plain infinitive \textbf{libban}. \textbf{Mid ðǣm smiðe} names the accompanying person in the dative. The intended sense is “he is unwilling to live with the blacksmith,” not a newly taught future auxiliary. A finite living verb in place of \textbf{libban} would break the volitional complement.
% mdparagraph:P122:end

\end{minipage}\par\addvspace{4pt}

\par\begin{minipage}{\linewidth}\raggedright
% item: EX-II-09-14
% key: KEY-EX-II-09-14
% mdparagraph:P123:begin
\textbf{Cook II.9.14.} We had rather live [To us it is preferable to sojourn] with the farmer than with the tailor.
% mdparagraph:P123:end

% mdparagraph:P124:begin
\textbf{Ūs is lēofre wician mid ðǣm ierðlinge ðonne mid ðǣm sēamere.}
% mdparagraph:P124:end

% mdparagraph:P125:begin
\textbf{Target and interpretation:} Dative experiencer \textbf{ūs}, neuter comparative \textbf{lēofre}, singular \textbf{is} and plain infinitive \textbf{wician} realize Cook's bracketed frame. The alternatives are parallel \textbf{mid + dative} phrases. English “had rather” does not assert past time here; \textbf{wæs} would change the preference's time. Natural alternative wording must still satisfy the stated frame.
% mdparagraph:P125:end

\end{minipage}\par\addvspace{4pt}

\par\begin{minipage}{\linewidth}\raggedright
% item: EX-II-09-15
% key: KEY-EX-II-09-15
% mdparagraph:P126:begin
\textbf{Cook II.9.15.} Gold is better (66) than tin.
% mdparagraph:P126:end

% mdparagraph:P127:begin
\textbf{Gold is bętere ðonne tin.}
% mdparagraph:P127:end

% mdparagraph:P128:begin
\textbf{Target and interpretation:} Neuter singular \textbf{gold} takes comparative \textbf{bętere}; nominative \textbf{tin} supplies the parallel standard after \textbf{ðonne}. Keep Cook's short \textbf{i} in \textbf{tin}. Masculine \textbf{bętera} would not agree with the gold; plain \textbf{gōd} would lose the comparison. The printed 66 is a grammar reference, not translation content.
% mdparagraph:P128:end

\end{minipage}\par\addvspace{4pt}

\par\begin{minipage}{\linewidth}\raggedright
% item: EX-II-09-16
% key: KEY-EX-II-09-16
% mdparagraph:P129:begin
\textbf{Cook II.9.16.} The carpenter makes a hole.
% mdparagraph:P129:end

% mdparagraph:P130:begin
\textbf{Se trēowwyrhta gewyrcð ðyrel.}
% mdparagraph:P130:end

% mdparagraph:P131:begin
\textbf{Target and interpretation:} The singular carpenter takes \textbf{gewyrcð}, with accusative neuter singular \textbf{ðyrel} as the thing made. The indefinite English article does not require inserting cardinal \textbf{ān}. The noun retains its medial vowel; a fabricated reduced stem would miss its lexical diagnostics.
% mdparagraph:P131:end

\end{minipage}\par\addvspace{4pt}

\reserve{7\baselineskip}
% mdparagraph:P132:begin
\section*{Cook Group II, Exercise 10}
% mdparagraph:P132:end

\par\begin{minipage}{\linewidth}\raggedright
% item: EX-II-10-01
% key: KEY-EX-II-10-01
% mdparagraph:P133:begin
\textbf{Cook II.10.1.} The tree is useful to the carpenter.
% mdparagraph:P133:end

% mdparagraph:P134:begin
\textbf{Ðæt trēo is behēfe ðǣm trēowwyrhtan.}
% mdparagraph:P134:end

% mdparagraph:P135:begin
\textbf{Target and interpretation:} Neuter subject \textbf{ðæt trēo} takes agreeing predicate \textbf{behēfe}, with dative singular recipient \textbf{ðǣm trēowwyrhtan}. The tree is useful to the carpenter; the carpenter is neither subject nor an accusative object. This lexical dative frame is explicitly taught, rather than inferred from the English adjective.
% mdparagraph:P135:end

\end{minipage}\par\addvspace{4pt}

\par\begin{minipage}{\linewidth}\raggedright
% item: EX-II-10-02
% key: KEY-EX-II-10-02
% mdparagraph:P136:begin
\textbf{Cook II.10.2.} Our horses have fodder.
% mdparagraph:P136:end

% mdparagraph:P137:begin
\textbf{Ūru hors habbað fōdor.}
% mdparagraph:P137:end

% mdparagraph:P138:begin
\textbf{Target and interpretation:} Nominative neuter plural possessive \textbf{ūru} agrees with plural \textbf{hors}, followed by plural \textbf{habbað}. \textbf{Fōdor} is the accusative singular mass noun fodder. The horses' plural has no added ending; do not infer singular from \textbf{hors} or make the verb singular. The possessive means our, not their.
% mdparagraph:P138:end

\end{minipage}\par\addvspace{4pt}

\par\begin{minipage}{\linewidth}\raggedright
% item: EX-II-10-03
% key: KEY-EX-II-10-03
% mdparagraph:P139:begin
\textbf{Cook II.10.3.} O priest, read the Gospel.
% mdparagraph:P139:end

% mdparagraph:P140:begin
\textbf{Ēalā mæsseprēost, rǣd ðæt godspell.}
% mdparagraph:P140:end

% mdparagraph:P141:begin
\textbf{Target and interpretation:} Nominative-shaped singular \textbf{mæsseprēost} is direct address; \textbf{rǣd} commands that one person to read accusative \textbf{ðæt godspell}. \textbf{Ēalā} renders “O.” Indicative \textbf{rǣtt} would state that someone reads; plural \textbf{rǣdað} would address several people instead of the named singular priest.
% mdparagraph:P141:end

\end{minipage}\par\addvspace{4pt}

\par\begin{minipage}{\linewidth}\raggedright
% item: EX-II-10-04
% key: KEY-EX-II-10-04
% mdparagraph:P142:begin
\textbf{Cook II.10.4.} Buy (sing.) wine and oil.
% mdparagraph:P142:end

% mdparagraph:P143:begin
\textbf{Byge wīn and ęle.}
% mdparagraph:P143:end

% mdparagraph:P144:begin
\textbf{Target and interpretation:} Expressly singular command \textbf{byge} takes two accusative singular material nouns, \textbf{wīn} and \textbf{ęle}. Two objects do not require plural command \textbf{bycgað}. The meaning is buy; substituting a selling verb would change the transaction.
% mdparagraph:P144:end

\end{minipage}\par\addvspace{4pt}

\par\begin{minipage}{\linewidth}\raggedright
% item: EX-II-10-05
% key: KEY-EX-II-10-05
% mdparagraph:P145:begin
\textbf{Cook II.10.5.} O good fellows, row over the waters.
% mdparagraph:P145:end

% mdparagraph:P146:begin
\textbf{Ēalā gōdan gefēran, rōwað ofer ðā wæteru.}
% mdparagraph:P146:end

% mdparagraph:P147:begin
\textbf{Target and interpretation:} Direct-address plural \textbf{gōdan gefēran} has a weak adjective and takes plural command \textbf{rōwað}. \textbf{Ofer ðā wæteru} gives the route over the waters, using accusative neuter plural in this model. The fellows are addressed, not objects; strong \textbf{gōde} would miss the selected weak-adjective address pattern.
% mdparagraph:P147:end

\end{minipage}\par\addvspace{4pt}

\par\begin{minipage}{\linewidth}\raggedright
% item: EX-II-10-06
% key: KEY-EX-II-10-06
% mdparagraph:P148:begin
\textbf{Cook II.10.6.} Her son is rich.
% mdparagraph:P148:end

% mdparagraph:P149:begin
\textbf{Hiere sunu is welig.}
% mdparagraph:P149:end

% mdparagraph:P150:begin
\textbf{Target and interpretation:} Invariant \textbf{hiere} means her. Nominative singular \textbf{sunu} takes singular \textbf{is} and masculine predicate \textbf{welig}. The possessor's gender does not make the son or predicate feminine. Genitive/\allowbreak{}dative \textbf{suna} would miss the son's subject case.
% mdparagraph:P150:end

\end{minipage}\par\addvspace{4pt}

\par\begin{minipage}{\linewidth}\raggedright
% item: EX-II-10-07
% key: KEY-EX-II-10-07
% mdparagraph:P151:begin
\textbf{Cook II.10.7.} Ply (sing.) your trade.
% mdparagraph:P151:end

% mdparagraph:P152:begin
\textbf{Begā ðīnne cræft.}
% mdparagraph:P152:end

% mdparagraph:P153:begin
\textbf{Target and interpretation:} Singular command \textbf{begā} takes accusative masculine singular \textbf{ðīnne cræft}, your trade. \textbf{Ðīnne} agrees with the possessed noun while retaining reference to one addressee. \textbf{Begāð} misses Cook's explicit singular command; \textbf{ðīn} is not the selected masculine accusative form.
% mdparagraph:P153:end

\end{minipage}\par\addvspace{4pt}

\par\begin{minipage}{\linewidth}\raggedright
% item: EX-II-10-08
% key: KEY-EX-II-10-08
% mdparagraph:P154:begin
\textbf{Cook II.10.8.} God gave us his Gospel.
% mdparagraph:P154:end

% mdparagraph:P155:begin
\textbf{God sealde ūs his godspell.}
% mdparagraph:P155:end

% mdparagraph:P156:begin
\textbf{Target and interpretation:} Preterite singular \textbf{sealde} has dative recipient \textbf{ūs} and accusative object \textbf{his godspell}. God gives the Gospel; we receive it. Invariant \textbf{his} refers to God. Present \textbf{sęlð} changes the time, and a genitive pronoun would not supply the recipient relation.
% mdparagraph:P156:end

\end{minipage}\par\addvspace{4pt}

\par\begin{minipage}{\linewidth}\raggedright
% item: EX-II-10-09
% key: KEY-EX-II-10-09
% mdparagraph:P157:begin
\textbf{Cook II.10.9.} The farmer practises agriculture.
% mdparagraph:P157:end

% mdparagraph:P158:begin
\textbf{Se ierðling begǣð eorðtilðe.}
% mdparagraph:P158:end

% mdparagraph:P159:begin
\textbf{Target and interpretation:} Singular \textbf{begǣð} takes accusative feminine singular \textbf{eorðtilðe}, the occupation practised. Nominative \textbf{eorðtilð} appears later as an answer naming agriculture, but is not this object form. The farmer engages in agriculture; he does not travel physically toward it.
% mdparagraph:P159:end

\end{minipage}\par\addvspace{4pt}

\par\begin{minipage}{\linewidth}\raggedright
% item: EX-II-10-10
% key: KEY-EX-II-10-10
% mdparagraph:P160:begin
\textbf{Cook II.10.10.} The service of God [God’s service] is better (66) than other occupations.
% mdparagraph:P160:end

% mdparagraph:P161:begin
\textbf{Godes ðēowdōm is bętera ðonne ōðre cræftas.}
% mdparagraph:P161:end

% mdparagraph:P162:begin
\textbf{Target and interpretation:} \textbf{Godes ðēowdōm} means service directed to God, an explicitly taught objective-genitive relationship. Masculine singular \textbf{ðēowdōm} takes \textbf{bętera}; nominative plural \textbf{ōðre cræftas} supplies the comparison after \textbf{ðonne}. Do not make the predicate plural from the comparison noun, or translate the service as work performed by God.
% mdparagraph:P162:end

\end{minipage}\par\addvspace{4pt}

\par\begin{minipage}{\linewidth}\raggedright
% item: EX-II-10-11
% key: KEY-EX-II-10-11
% mdparagraph:P163:begin
\textbf{Cook II.10.11.} Which of you holds (om. the) chief place?
% mdparagraph:P163:end

% mdparagraph:P164:begin
\textbf{Hwilc ēower hielt ealdorscipe?}
% mdparagraph:P164:end

% mdparagraph:P165:begin
\textbf{Target and interpretation:} Substantival nominative \textbf{hwilc} selects one member of the group named by personal genitive plural \textbf{ēower}, which of you. Singular \textbf{hielt} takes accusative \textbf{ealdorscipe}, without an article as requested. Here \textbf{ēower} is a personal genitive, not an agreeing possessive. Retain the selected short-vowel spelling \textbf{hielt}.
% mdparagraph:P165:end

\end{minipage}\par\addvspace{4pt}

\par\begin{minipage}{\linewidth}\raggedright
% item: EX-II-10-12
% key: KEY-EX-II-10-12
% mdparagraph:P166:begin
\textbf{Cook II.10.12.} My gain is little.
% mdparagraph:P166:end

% mdparagraph:P167:begin
\textbf{Mīn gestrēon is lȳtel.}
% mdparagraph:P167:end

% mdparagraph:P168:begin
\textbf{Target and interpretation:} Nominative neuter singular \textbf{mīn gestrēon} takes strong neuter predicate \textbf{lȳtel}. The gain is little or small; comparative less would change the assertion. The possessive agrees with the singular gain. A weak form copied from a definite attributive phrase would miss the predicate construction.
% mdparagraph:P168:end

\end{minipage}\par\addvspace{4pt}

\par\begin{minipage}{\linewidth}\raggedright
% item: EX-II-10-13
% key: KEY-EX-II-10-13
% mdparagraph:P169:begin
\textbf{Cook II.10.13.} My labor is great.
% mdparagraph:P169:end

% mdparagraph:P170:begin
\textbf{Mīn gedeorf is micel.}
% mdparagraph:P170:end

% mdparagraph:P171:begin
\textbf{Target and interpretation:} Nominative neuter singular \textbf{mīn gedeorf} takes \textbf{is} and strong predicate \textbf{micel}. The statement concerns the magnitude of my labor, not a comparison or the number of workers. Keep the possessive's reference to the speaker.
% mdparagraph:P171:end

\end{minipage}\par\addvspace{4pt}

\par\begin{minipage}{\linewidth}\raggedright
% item: EX-II-10-14
% key: KEY-EX-II-10-14
% mdparagraph:P172:begin
\textbf{Cook II.10.14.} I sold ivory and silk.
% mdparagraph:P172:end

% mdparagraph:P173:begin
\textbf{Ic cīepte elpes bān and sīdan.}
% mdparagraph:P173:end

% mdparagraph:P174:begin
\textbf{Target and interpretation:} First-singular preterite \textbf{cīepte} means sold. Supplied \textbf{elpes bān} has genitive dependent \textbf{elpes} and accusative head \textbf{bān}; weak feminine \textbf{sīde} supplies accusative \textbf{sīdan}. \textbf{Bohte} would mean bought and reverse the transaction. The complete ivory expression stays available; no general rule for inventing material compounds is tested.
% mdparagraph:P174:end

\end{minipage}\par\addvspace{4pt}

\par\begin{minipage}{\linewidth}\raggedright
% item: EX-II-10-15
% key: KEY-EX-II-10-15
% mdparagraph:P175:begin
\textbf{Cook II.10.15.} O, were I king of this people!
% mdparagraph:P175:end

% mdparagraph:P176:begin
\textbf{Ēalā, wǣre ic cyning ðisses folces!}
% mdparagraph:P176:end

% mdparagraph:P177:begin
\textbf{Target and interpretation:} First-singular preterite subjunctive \textbf{wǣre} expresses an unreal present wish, with nominative predicate \textbf{cyning} and genitive neuter singular \textbf{ðisses folces}. It does not assert that the speaker formerly was king. Optional \textbf{ēalā} reinforces the wish. The predicate's case follows its role, even though \textbf{cyning} has identical nominative and accusative spellings.
% mdparagraph:P177:end

\end{minipage}\par\addvspace{4pt}

\reserve{7\baselineskip}
% mdparagraph:P178:begin
\section*{Cook Group II, Exercise 11}
% mdparagraph:P178:end

\par\begin{minipage}{\linewidth}\raggedright
% item: EX-II-11-01
% key: KEY-EX-II-11-01
% mdparagraph:P179:begin
\textbf{Cook II.11.1.} The monks have food and drink.
% mdparagraph:P179:end

% mdparagraph:P180:begin
\textbf{Ðā munucas habbað męte and dręnc.}
% mdparagraph:P180:end

% mdparagraph:P181:begin
\textbf{Target and interpretation:} Nominative plural \textbf{ðā munucas} takes \textbf{habbað}, with accusative singular mass nouns \textbf{męte/\allowbreak{}dręnc}, food and drink. The number of objects does not determine the finite agreement. One monk and a singular verb would change the specified participants.
% mdparagraph:P181:end

\end{minipage}\par\addvspace{4pt}

\par\begin{minipage}{\linewidth}\raggedright
% item: EX-II-11-02
% key: KEY-EX-II-11-02
% mdparagraph:P182:begin
\textbf{Cook II.11.2.} Settle [quell] your disputes.
% mdparagraph:P182:end

% mdparagraph:P183:begin
\textbf{Gestillað ēowru geflīt.}
% mdparagraph:P183:end

% mdparagraph:P184:begin
\textbf{Target and interpretation:} The model addresses several people: plural \textbf{gestillað} and accusative neuter plural \textbf{ēowru geflīt}, your disputes. English leaves addressee number open; \textbf{Gestil ðīn geflīt} is valid for one person. Match command and possessive consistently. Endingless \textbf{geflīt} is plural here; settle has the hinted sense quell, not establish a residence.
% mdparagraph:P184:end

\end{minipage}\par\addvspace{4pt}

\par\begin{minipage}{\linewidth}\raggedright
% item: EX-II-11-03
% key: KEY-EX-II-11-03
% mdparagraph:P185:begin
\textbf{Cook II.11.3.} He suffered shipwreck.
% mdparagraph:P185:end

% mdparagraph:P186:begin
\textbf{Hē ðolode forlidennesse.}
% mdparagraph:P186:end

% mdparagraph:P187:begin
\textbf{Target and interpretation:} Third-singular preterite \textbf{ðolode} takes accusative feminine singular \textbf{forlidennesse}. The \textbf{-ness-} stem appears before the noun ending. He suffered shipwreck; the sentence does not say he caused it or died in it. Present \textbf{ðolað} would change the time.
% mdparagraph:P187:end

\end{minipage}\par\addvspace{4pt}

\par\begin{minipage}{\linewidth}\raggedright
% item: EX-II-11-04
% key: KEY-EX-II-11-04
% mdparagraph:P188:begin
\textbf{Cook II.11.4.} He abandoned his wife.
% mdparagraph:P188:end

% mdparagraph:P189:begin
\textbf{Hē forlēt his wīf.}
% mdparagraph:P189:end

% mdparagraph:P190:begin
\textbf{Target and interpretation:} Singular preterite \textbf{forlēt} takes accusative singular \textbf{his wīf}, the wife abandoned. \textbf{His} refers to the male subject; \textbf{wīf} is grammatically neuter despite referring to a woman. Present \textbf{forlǣtt} would describe a present abandonment rather than the stated past event.
% mdparagraph:P190:end

\end{minipage}\par\addvspace{4pt}

\par\begin{minipage}{\linewidth}\raggedright
% item: EX-II-11-05
% key: KEY-EX-II-11-05
% mdparagraph:P191:begin
\textbf{Cook II.11.5.} Be what [that] you (sing.) ought to be.
% mdparagraph:P191:end

% mdparagraph:P192:begin
\textbf{Bēo ðæt ðe ðū scealt wesan.}
% mdparagraph:P192:end

% mdparagraph:P193:begin
\textbf{Target and interpretation:} Singular command \textbf{bēo} is followed by bounded predicative \textbf{ðæt ðe}, what or that which. Inside the relative, \textbf{ðū} takes \textbf{scealt} with plain \textbf{wesan}. Relative \textbf{ðæt} is a nominative predicate, not an object of being or an invariant content conjunction. Prior \textbf{wes} is a defensible singular command alternative; plural \textbf{bēoð} misses the specified number.
% mdparagraph:P193:end

\end{minipage}\par\addvspace{4pt}

\par\begin{minipage}{\linewidth}\raggedright
% item: EX-II-11-06
% key: KEY-EX-II-11-06
% mdparagraph:P194:begin
\textbf{Cook II.11.6.} They do not wish to be what they ought to be.
% mdparagraph:P194:end

% mdparagraph:P195:begin
\textbf{Hīe ne willað wesan ðæt ðe hīe sculon wesan.}
% mdparagraph:P195:end

% mdparagraph:P196:begin
\textbf{Target and interpretation:} Plural \textbf{hīe ne willað} expresses unwillingness, followed by \textbf{wesan} and the bounded predicative relative. Within that relative, plural \textbf{hīe} takes \textbf{sculon} with another \textbf{wesan}. One infinitive completes wishing, the other obligation. Neither takes plural finite endings. Preserve obligation, rather than converting it into a simple future prediction.
% mdparagraph:P196:end

\end{minipage}\par\addvspace{4pt}

\par\begin{minipage}{\linewidth}\raggedright
% item: EX-II-11-07
% key: KEY-EX-II-11-07
% mdparagraph:P197:begin
\textbf{Cook II.11.7.} He went aboard the ship with his wares.
% mdparagraph:P197:end

% mdparagraph:P198:begin
\textbf{Hē āstāg ðæt scip mid his hlæstum.}
% mdparagraph:P198:end

% mdparagraph:P199:begin
\textbf{Target and interpretation:} Preterite singular \textbf{āstāg} takes direct accusative ship \textbf{ðæt scip}; dative \textbf{mid his hlæstum} names accompanying wares. Documented \textbf{āstāh} is a defensible past-singular variant. Do not make the ship's case follow \textbf{mid}, which governs the wares, or replace the direct-object frame with an unexplained destination preposition.
% mdparagraph:P199:end

\end{minipage}\par\addvspace{4pt}

\par\begin{minipage}{\linewidth}\raggedright
% item: EX-II-11-08
% key: KEY-EX-II-11-08
% mdparagraph:P200:begin
\textbf{Cook II.11.8.} O good workmen, ply your trades.
% mdparagraph:P200:end

% mdparagraph:P201:begin
\textbf{Ēalā gōdan wyrhtan, begāð ēowre cræftas.}
% mdparagraph:P201:end

% mdparagraph:P202:begin
\textbf{Target and interpretation:} Weak plural \textbf{gōdan} accompanies nominative-shaped addressed \textbf{wyrhtan}, followed by plural \textbf{begāð}. \textbf{Ēowre cræftas} is accusative masculine plural, your trades. Singular \textbf{begā} would not address the plural workmen as a group. The address adjective and possessive have different environments and endings.
% mdparagraph:P202:end

\end{minipage}\par\addvspace{4pt}

\par\begin{minipage}{\linewidth}\raggedright
% item: EX-II-11-09
% key: KEY-EX-II-11-09
% mdparagraph:P203:begin
\textbf{Cook II.11.9.} I tell you (sing.), the shoemaker buys awls, and the fisherman buys (fish-) hooks.
% mdparagraph:P203:end

% mdparagraph:P204:begin
\textbf{Ic sęcge ðē, se scēowyrhta bygeð ǣla and se fiscere bygeð anglas.}
% mdparagraph:P204:end

% mdparagraph:P205:begin
\textbf{Target and interpretation:} First-singular \textbf{sęcge} takes dative singular \textbf{ðē}; the direct/\allowbreak{}paratactic statement then has two singular \textbf{bygeð} clauses. \textbf{Ǣla/\allowbreak{}anglas} are the goods bought. Cook prints no “that,” and this route needs none. A clearly signaled conversion to an indirect report can be valid but is not forced; substituting selling would change the meaning.
% mdparagraph:P205:end

\end{minipage}\par\addvspace{4pt}

\par\begin{minipage}{\linewidth}\raggedright
% item: EX-II-11-10
% key: KEY-EX-II-11-10
% mdparagraph:P206:begin
\textbf{Cook II.11.10.} Sometimes he reads, because he is wise.
% mdparagraph:P206:end

% mdparagraph:P207:begin
\textbf{Hwīlum hē rǣtt, for-ðām hē wīs is.}
% mdparagraph:P207:end

% mdparagraph:P208:begin
\textbf{Target and interpretation:} Frequency adverb \textbf{hwīlum} modifies present singular \textbf{rǣtt}; \textbf{for-ðām hē wīs is} gives the factual reason. Wisdom is the reason, not the result. \textbf{Wīs} agrees with singular \textbf{hē}, and indicative \textbf{is} suits the asserted cause. The conjunction does not automatically demand subjunctive; sometimes does not mean always.
% mdparagraph:P208:end

\end{minipage}\par\addvspace{4pt}

\par\begin{minipage}{\linewidth}\raggedright
% item: EX-II-11-11
% key: KEY-EX-II-11-11
% mdparagraph:P209:begin
\textbf{Cook II.11.11.} The king buys horses.
% mdparagraph:P209:end

% mdparagraph:P210:begin
\textbf{Se cyning bygeð hors.}
% mdparagraph:P210:end

% mdparagraph:P211:begin
\textbf{Target and interpretation:} Singular \textbf{se cyning bygeð} takes accusative plural \textbf{hors}, the horses bought. Heavy neuter \textbf{hors} has identical singular and plural nominative/\allowbreak{}accusative spellings; the English task determines plural here. The king is one person, so the object's plurality does not call for \textbf{bycgað} as his verb.
% mdparagraph:P211:end

\end{minipage}\par\addvspace{4pt}

\par\begin{minipage}{\linewidth}\raggedright
% item: EX-II-11-12
% key: KEY-EX-II-11-12
% mdparagraph:P212:begin
\textbf{Cook II.11.12.} Righteousness is better than gold.
% mdparagraph:P212:end

% mdparagraph:P213:begin
\textbf{Rihtwīsnes is bętere ðonne gold.}
% mdparagraph:P213:end

% mdparagraph:P214:begin
\textbf{Target and interpretation:} Feminine subject \textbf{rihtwīsnes} takes weak comparative nominative feminine singular \textbf{bętere}. Neuter \textbf{gold} is the parallel nominative standard after \textbf{ðonne}. Predicate agreement follows righteousness, not gold; masculine \textbf{bętera} misses that agreement. The sentence compares value rather than identifying the two things.
% mdparagraph:P214:end

\end{minipage}\par\addvspace{4pt}

\par\begin{minipage}{\linewidth}\raggedright
% item: EX-II-11-13
% key: KEY-EX-II-11-13
% mdparagraph:P215:begin
\textbf{Cook II.11.13.} They struck me.
% mdparagraph:P215:end

% mdparagraph:P216:begin
\textbf{Hīe slōgon mē.}
% mdparagraph:P216:end

% mdparagraph:P217:begin
\textbf{Target and interpretation:} Nominative plural \textbf{hīe} takes preterite plural \textbf{slōgon}; \textbf{mē} is the accusative person struck. This is the object-taking sense of \textbf{slēan}. Singular \textbf{slōh} loses the plural subject. Syncretic accusative/\allowbreak{}dative \textbf{mē} does not make the object's role dative under this verb.
% mdparagraph:P217:end

\end{minipage}\par\addvspace{4pt}

\par\begin{minipage}{\linewidth}\raggedright
% item: EX-II-11-14
% key: KEY-EX-II-11-14
% mdparagraph:P218:begin
\textbf{Cook II.11.14.} The farmer beats him.
% mdparagraph:P218:end

% mdparagraph:P219:begin
\textbf{Se ierðling bīett hine.}
% mdparagraph:P219:end

% mdparagraph:P220:begin
\textbf{Target and interpretation:} The singular farmer takes present \textbf{bīett} and accusative masculine singular \textbf{hine}, the person beaten. Keep the selected present stem mutation and dental contraction. Plural \textbf{bēatað} does not agree with this subject; nominative \textbf{hē} does not supply the object.
% mdparagraph:P220:end

\end{minipage}\par\addvspace{4pt}

\par\begin{minipage}{\linewidth}\raggedright
% item: EX-II-11-15
% key: KEY-EX-II-11-15
% mdparagraph:P221:begin
\textbf{Cook II.11.15.} We suffer the loss of copper and tin.
% mdparagraph:P221:end

% mdparagraph:P222:begin
\textbf{Wē ðoliað ǣres and tines lyre.}
% mdparagraph:P222:end

% mdparagraph:P223:begin
\textbf{Target and interpretation:} Present plural \textbf{ðoliað} takes accusative singular \textbf{lyre}, the loss suffered. Genitives \textbf{ǣres/\allowbreak{}tines} identify the materials lost and precede the head. The speakers suffer loss of those metals, not simply the metals themselves. Preserve the long vowel in copper \textbf{ǣr} but Cook's short \textbf{i} in \textbf{tin/\allowbreak{}tines}.
% mdparagraph:P223:end

\end{minipage}\par\addvspace{4pt}

\par\begin{minipage}{\linewidth}\raggedright
% item: EX-II-11-16
% key: KEY-EX-II-11-16
% mdparagraph:P224:begin
\textbf{Cook II.11.16.} The princes of the earth go aboard their ships.
% mdparagraph:P224:end

% mdparagraph:P225:begin
\textbf{Eorðan ealdormęnn āstīgað hiera scipu.}
% mdparagraph:P225:end

% mdparagraph:P226:begin
\textbf{Target and interpretation:} Genitive \textbf{eorðan} depends on nominative plural \textbf{ealdormęnn}, governing plural \textbf{āstīgað}. \textbf{Hiera scipu} is accusative plural, their ships, with invariant possessor \textbf{hiera}. Explicit \textbf{ðǣre} with genitive \textbf{eorðan} is also defensible. The princes' plural uses \textbf{mann/\allowbreak{}męnn} in the final member, not a mechanically added \textbf{-as}.
% mdparagraph:P226:end

\end{minipage}\par\addvspace{4pt}

\par\begin{minipage}{\linewidth}\raggedright
% item: EX-II-11-17
% key: KEY-EX-II-11-17
% mdparagraph:P227:begin
\textbf{Cook II.11.17.} The priest benefits (164. e) the laymen.
% mdparagraph:P227:end

% mdparagraph:P228:begin
\textbf{Se mæsseprēost fręmeð ðǣm ceorlum.}
% mdparagraph:P228:end

% mdparagraph:P229:begin
\textbf{Target and interpretation:} Singular \textbf{fręmeð} in the benefit sense governs dative plural \textbf{ðǣm ceorlum}. The priest is subject; the laymen are beneficiaries. Accusative \textbf{ðā ceorlas} would miss this lexical government. Cook's \textbf{164. e} refers to the dative rule and is not additional content to translate.
% mdparagraph:P229:end

\end{minipage}\par\addvspace{4pt}

\reserve{7\baselineskip}
% mdparagraph:P230:begin
\section*{Cook Group II, Exercise 12}
% mdparagraph:P230:end

\par\begin{minipage}{\linewidth}\raggedright
% item: EX-II-12-01
% key: KEY-EX-II-12-01
% mdparagraph:P231:begin
\textbf{Cook II.12.1.} The counsellor is wiser (64) than the warrior.
% mdparagraph:P231:end

% mdparagraph:P232:begin
\textbf{Se geðeahtend is wīsra ðonne se beorn.}
% mdparagraph:P232:end

% mdparagraph:P233:begin
\textbf{Target and interpretation:} Singular \textbf{se geðeahtend} takes weak comparative nominative masculine singular \textbf{wīsra}. Nominative \textbf{se beorn} supplies the parallel standard after \textbf{ðonne}. The adjective means wiser, not merely wise. The new \textbf{-end} agent is not an agreeing verbal participle; the warrior is not an accusative object.
% mdparagraph:P233:end

\end{minipage}\par\addvspace{4pt}

\par\begin{minipage}{\linewidth}\raggedright
% item: EX-II-12-02
% key: KEY-EX-II-12-02
% mdparagraph:P234:begin
\textbf{Cook II.12.2.} The elephant is larger than the horse.
% mdparagraph:P234:end

% mdparagraph:P235:begin
\textbf{Se elp is māra ðonne ðæt hors.}
% mdparagraph:P235:end

% mdparagraph:P236:begin
\textbf{Target and interpretation:} Masculine singular \textbf{se elp} takes weak comparative \textbf{māra}. \textbf{Ðæt hors} is the parallel nominative neuter singular comparison. The comparative agrees with the elephant, not the horse. \textbf{Micel} would express large without retaining larger.
% mdparagraph:P236:end

\end{minipage}\par\addvspace{4pt}

\par\begin{minipage}{\linewidth}\raggedright
% item: EX-II-12-03
% key: KEY-EX-II-12-03
% mdparagraph:P237:begin
\textbf{Cook II.12.3.} God created all (om. the) fishes in (om. the) sea.
% mdparagraph:P237:end

% mdparagraph:P238:begin
\textbf{God scēop ealle fiscas on sǣ.}
% mdparagraph:P238:end

% mdparagraph:P239:begin
\textbf{Target and interpretation:} Preterite singular \textbf{scēop} takes accusative masculine plural \textbf{ealle fiscas}, with location \textbf{on sǣ}. Both article omissions follow Cook's controls. Neuter plural \textbf{eall} would not agree with masculine fishes. Prior \textbf{gescēop} is also defensible creation vocabulary if the time and case relations are retained.
% mdparagraph:P239:end

\end{minipage}\par\addvspace{4pt}

\par\begin{minipage}{\linewidth}\raggedright
% item: EX-II-12-04
% key: KEY-EX-II-12-04
% mdparagraph:P240:begin
\textbf{Cook II.12.4.} The cattle seek fodder.
% mdparagraph:P240:end

% mdparagraph:P241:begin
\textbf{Ðā nīetenu sēcað fōdor.}
% mdparagraph:P241:end

% mdparagraph:P242:begin
\textbf{Target and interpretation:} Nominative neuter plural \textbf{ðā nīetenu} takes present plural indicative \textbf{sēcað}, with accusative singular mass noun \textbf{fōdor}. The cattle seek; the fodder is sought. \textbf{Sēcað} shares a spelling with the plural command, but the expressed third-person subject and statement make it indicative here.
% mdparagraph:P242:end

\end{minipage}\par\addvspace{4pt}

\par\begin{minipage}{\linewidth}\raggedright
% item: EX-II-12-05
% key: KEY-EX-II-12-05
% mdparagraph:P243:begin
\textbf{Cook II.12.5.} God commanded the earth to bring forth grass.
% mdparagraph:P243:end

% mdparagraph:P244:begin
\textbf{God hēt ðā eorðan gærs tēon forð.}
% mdparagraph:P244:end

% mdparagraph:P245:begin
\textbf{Target and interpretation:} Past command \textbf{hēt} takes accusative feminine singular actor \textbf{ðā eorðan} and plain \textbf{tēon} with \textbf{forð}, bring forth. \textbf{Gærs} is the infinitive's accusative object. Earth and grass are both accusative but have different functions: earth produces, grass is produced. An independent finite wish would miss the reported-command frame.
% mdparagraph:P245:end

\end{minipage}\par\addvspace{4pt}

\par\begin{minipage}{\linewidth}\raggedright
% item: EX-II-12-06
% key: KEY-EX-II-12-06
% mdparagraph:P246:begin
\textbf{Cook II.12.6.} On the third day God gathered the waters that were under the heavens.
% mdparagraph:P246:end

% mdparagraph:P247:begin
\textbf{On ðone ðriddan dæg gegaderode God ðā wæteru ðe under ðǣm heofonum wǣron.}
% mdparagraph:P247:end

% mdparagraph:P248:begin
\textbf{Target and interpretation:} Temporal \textbf{on ðone ðriddan dæg} is accusative masculine singular with weak ordinal. Singular \textbf{gegaderode} agrees with \textbf{God}; \textbf{ðā wæteru} is the accusative object. Within the relative, \textbf{ðe} is subject of plural \textbf{wǣron}, with dative location \textbf{under ðǣm heofonum}. Do not make the relative verb singular from main subject \textbf{God}.
% mdparagraph:P248:end

\end{minipage}\par\addvspace{4pt}

\par\begin{minipage}{\linewidth}\raggedright
% item: EX-II-12-07
% key: KEY-EX-II-12-07
% mdparagraph:P249:begin
\textbf{Cook II.12.7.} Let the earth bring forth living creatures (the verb precedes).
% mdparagraph:P249:end

% mdparagraph:P250:begin
\textbf{Tēo sēo eorðe libbende gesceafta forð.}
% mdparagraph:P250:end

% mdparagraph:P251:begin
\textbf{Target and interpretation:} Initial present subjunctive singular \textbf{tēo} expresses the third-person wish, followed by nominative \textbf{sēo eorðe}, meeting “the verb precedes.” Supplied \textbf{libbende gesceafta} is feminine accusative plural, using the documented feminine \textbf{-e} participial variant. Prior neuter \textbf{libbendu} does not agree. No English-like let verb or independent recall of the supplied participle is required.
% mdparagraph:P251:end

\end{minipage}\par\addvspace{4pt}

\par\begin{minipage}{\linewidth}\raggedright
% item: EX-II-12-08
% key: KEY-EX-II-12-08
% mdparagraph:P252:begin
\textbf{Cook II.12.8.} God made man good (173).
% mdparagraph:P252:end

% mdparagraph:P253:begin
\textbf{God geworhte ðone mann gōdne.}
% mdparagraph:P253:end

% mdparagraph:P254:begin
\textbf{Target and interpretation:} \textbf{Geworhte} takes accusative \textbf{ðone mann} and strong masculine accusative predicate \textbf{gōdne}. God makes; goodness is predicated of the man. Bare \textbf{mann} is also defensible for generic man. Nominative \textbf{gōd} or weak \textbf{gōdan} misses the selected object-predicate agreement; the article does not make the separate predicate weak.
% mdparagraph:P254:end

\end{minipage}\par\addvspace{4pt}

\par\begin{minipage}{\linewidth}\raggedright
% item: EX-II-12-09
% key: KEY-EX-II-12-09
% mdparagraph:P255:begin
\textbf{Cook II.12.9.} He is (om. the) light that gives light to (ālīehtan) all things.
% mdparagraph:P255:end

% mdparagraph:P256:begin
\textbf{Hē is lēoht ðe eall ðing ālīehteð.}
% mdparagraph:P256:end

% mdparagraph:P257:begin
\textbf{Target and interpretation:} Bare nominative predicate \textbf{lēoht} follows Cook's omission instruction. Relative \textbf{ðe} is subject of singular \textbf{ālīehteð}; \textbf{eall ðing} is accusative neuter plural, the things illuminated. English “gives light to” does not make them dative under \textbf{ālīehtan}. Retain the selected uncontracted present ending and the relative clause.
% mdparagraph:P257:end

\end{minipage}\par\addvspace{4pt}

\par\begin{minipage}{\linewidth}\raggedright
% item: EX-II-12-10
% key: KEY-EX-II-12-10
% mdparagraph:P258:begin
\textbf{Cook II.12.10.} Some of them (151) were rich, and some were wise.
% mdparagraph:P258:end

% mdparagraph:P259:begin
\textbf{Hīe sume wǣron welige, and sume wǣron wīse.}
% mdparagraph:P259:end

% mdparagraph:P260:begin
\textbf{Target and interpretation:} Selected apposition \textbf{hīe sume}, both nominative plural, means some of them. Second \textbf{sume} supplies the subject of \textbf{wǣron wīse}. Strong masculine plural \textbf{welige/\allowbreak{}wīse} agrees with the people. \textbf{Sume hiera} is a valid genitive-partitive alternative, but \textbf{hīe} itself is not genitive. No decision about overlap between the wealthy and wise groups is required.
% mdparagraph:P260:end

\end{minipage}\par\addvspace{4pt}

\par\begin{minipage}{\linewidth}\raggedright
% item: EX-II-12-11
% key: KEY-EX-II-12-11
% mdparagraph:P261:begin
\textbf{Cook II.12.11.} The night has many stars.
% mdparagraph:P261:end

% mdparagraph:P262:begin
\textbf{Sēo niht hæfð manige steorran.}
% mdparagraph:P262:end

% mdparagraph:P263:begin
\textbf{Target and interpretation:} Singular \textbf{sēo niht hæfð} takes accusative masculine plural \textbf{manige steorran}. \textbf{Manige} is ordinary strong attributive many; \textbf{steorran} is the weak noun's plural. The object does not make the night or verb plural. Singular many-a, substantival many and a new genitive quantity construction are not assigned.
% mdparagraph:P263:end

\end{minipage}\par\addvspace{4pt}

\par\begin{minipage}{\linewidth}\raggedright
% item: EX-II-12-12
% key: KEY-EX-II-12-12
% mdparagraph:P264:begin
\textbf{Cook II.12.12.} The light shineth in (on) darkness (plur.).
% mdparagraph:P264:end

% mdparagraph:P265:begin
\textbf{Ðæt lēoht scīnð on ðīestrum.}
% mdparagraph:P265:end

% mdparagraph:P266:begin
\textbf{Target and interpretation:} Singular \textbf{ðæt lēoht scīnð} has dative plural location \textbf{on ðīestrum}. The selected darkness noun is plural, as Cook reminds the learner; do not invent a singular to match English. Present \textbf{scīnð} retains long \textbf{ī}, unlike the short-vowel preterite plural/\allowbreak{}subjunctive stem.
% mdparagraph:P266:end

\end{minipage}\par\addvspace{4pt}

\par\begin{minipage}{\linewidth}\raggedright
% item: EX-II-12-13
% key: KEY-EX-II-12-13
% mdparagraph:P267:begin
\textbf{Cook II.12.13.} The year has three hundred and sixty-five (78) days.
% mdparagraph:P267:end

% mdparagraph:P268:begin
\textbf{Ðæt gēar hæfð ðrēo hund and fīf and siextig daga.}
% mdparagraph:P268:end

% mdparagraph:P269:begin
\textbf{Target and interpretation:} Singular \textbf{ðæt gēar hæfð} takes the complete supplied 365-days expression. Within the help, \textbf{ðrēo} counts neuter \textbf{hund}; \textbf{fīf and siextig} gives five and sixty; \textbf{daga} is genitive plural. This is a declared rule adaptation, not an exact quotation or a full higher-numeral assignment. Masculine \textbf{ðrīe} does not replace the supplied neuter form.
% mdparagraph:P269:end

\end{minipage}\par\addvspace{4pt}

\par\begin{minipage}{\linewidth}\raggedright
% item: EX-II-12-14
% key: KEY-EX-II-12-14
% mdparagraph:P270:begin
\textbf{Cook II.12.14.} The carpenter has two sons.
% mdparagraph:P270:end

% mdparagraph:P271:begin
\textbf{Se trēowwyrhta hæfð twēgen suna.}
% mdparagraph:P271:end

% mdparagraph:P272:begin
\textbf{Target and interpretation:} Singular carpenter and \textbf{hæfð} take accusative masculine plural \textbf{twēgen suna}. Selected \textbf{sunu} has plural \textbf{suna}, not an obligatory \textbf{-as} ending. Masculine \textbf{twēgen}, not neuter/\allowbreak{}feminine \textbf{twā}, agrees with sons. Documented later \textbf{sunas} can be valid language while \textbf{suna} remains the assigned u-stem form.
% mdparagraph:P272:end

\end{minipage}\par\addvspace{4pt}

\par\begin{minipage}{\linewidth}\raggedright
% item: EX-II-12-15
% key: KEY-EX-II-12-15
% mdparagraph:P273:begin
\textbf{Cook II.12.15.} The soldier has eight houses.
% mdparagraph:P273:end

% mdparagraph:P274:begin
\textbf{Se cęmpa hæfð eahta hūs.}
% mdparagraph:P274:end

% mdparagraph:P275:begin
\textbf{Target and interpretation:} Singular \textbf{se cęmpa hæfð} takes invariant preposed \textbf{eahta} and accusative neuter plural \textbf{hūs}. The heavy neuter plural adds no \textbf{-u}; its identical singular spelling does not erase the eight houses. The soldier is one person, so the verb stays singular.
% mdparagraph:P275:end

\end{minipage}\par\addvspace{4pt}

\par\begin{minipage}{\linewidth}\raggedright
% item: EX-II-12-16
% key: KEY-EX-II-12-16
% mdparagraph:P276:begin
\textbf{Cook II.12.16.} The peril was great.
% mdparagraph:P276:end

% mdparagraph:P277:begin
\textbf{Se pliht wæs micel.}
% mdparagraph:P277:end

% mdparagraph:P278:begin
\textbf{Target and interpretation:} Masculine singular \textbf{se pliht} takes preterite singular \textbf{wæs} and strong masculine predicate \textbf{micel}. The danger was great, not a present assertion or comparison. Weak \textbf{micla} would miss the selected ordinary predicate form.
% mdparagraph:P278:end

\end{minipage}\par\addvspace{4pt}

\reserve{7\baselineskip}
% mdparagraph:P279:begin
\section*{R. The fisherman and his companion}
% mdparagraph:P279:end

\par\begin{minipage}{\linewidth}\raggedright
% item: NEW-U09-R01
% key: KEY-NEW-U09-R01
% mdparagraph:P280:begin
\textbf{R01. Who asks, and who has food?}
% mdparagraph:P280:end

% mdparagraph:P281:begin
The companion requests information: \textbf{Se gefēra cwið: “Saga mē, hwæt ðū hæbbe.”} — The companion says, “Tell me what you have.”
% mdparagraph:P281:end

% mdparagraph:P282:begin
The fisherman is reported to have food: \textbf{Se fiscere sægð, ðæt hē męte hæbbe.} — The fisherman says that he has food.
% mdparagraph:P282:end

% mdparagraph:P283:begin
\textbf{Target and interpretation:} Identify the two speakers and retain \textbf{hē} as the fisherman. The companion's \textbf{mē} names the recipient of the requested information, not the possessor of the food. Saying the companion has the food at this point reverses the reported reference.
% mdparagraph:P283:end

\end{minipage}\par\addvspace{4pt}

\par\begin{minipage}{\linewidth}\raggedright
% item: NEW-U09-R02
% key: KEY-NEW-U09-R02
% mdparagraph:P284:begin
\textbf{R02. Three dependent clauses.}
% mdparagraph:P284:end

\begin{itemize}
% mdparagraph:P285:begin
\item \textbf{Hwæt ðū hæbbe} — what you have: a dependent question, with present subjunctive \textbf{hæbbe}.
% mdparagraph:P285:end
% mdparagraph:P286:begin
\item \textbf{Ðæt hē męte hæbbe} — that he has food: reported content, also with present subjunctive \textbf{hæbbe}.
% mdparagraph:P286:end
% mdparagraph:P287:begin
\item \textbf{Ðæt ðū mē hilpst} — that you help me: asserted factual content after knowing, with present indicative \textbf{hilpst}.
% mdparagraph:P287:end
\end{itemize}

% mdparagraph:P288:begin
\textbf{Target and interpretation:} Distinguish question, report and asserted knowledge. The two \textbf{hæbbe} forms are subjunctive here, but one has second-singular \textbf{ðū}, the other third-singular \textbf{hē}. Final \textbf{mē} is dative under \textbf{helpan}. Neither every \textbf{ðæt} clause nor every clause after knowing requires one mechanically prescribed mood.
% mdparagraph:P288:end

\end{minipage}\par\addvspace{4pt}

\par\begin{minipage}{\linewidth}\raggedright
% item: NEW-U09-R03
% key: KEY-NEW-U09-R03
% mdparagraph:P289:begin
\textbf{R03. The companion as actor.}
% mdparagraph:P289:end

% mdparagraph:P290:begin
In \textbf{Hē hǣtt his gefēran ðone fisc beran}, finite \textbf{hǣtt} commands accusative singular actor \textbf{his gefēran}, the fisherman's companion.
% mdparagraph:P290:end

% mdparagraph:P291:begin
In \textbf{Ic gesēo ðē ðone fisc beran}, finite \textbf{gesēo} perceives accusative singular actor \textbf{ðē}, the companion now addressed directly.
% mdparagraph:P291:end

% mdparagraph:P292:begin
In both, \textbf{ðone fisc} is the accusative singular object of \textbf{beran}, distinct from its actor.
% mdparagraph:P292:end

% mdparagraph:P293:begin
\textbf{Target and interpretation:} The carrier is the same companion, referred to in third person and then addressed in second person. \textbf{Beran} remains nonfinite. \textbf{Gefēran} has other possible case/\allowbreak{}number analyses elsewhere and \textbf{ðē} can elsewhere be dative; the command/\allowbreak{}perception frames settle their accusative actor roles here.
% mdparagraph:P293:end

\end{minipage}\par\addvspace{4pt}

\par\begin{minipage}{\linewidth}\raggedright
% item: NEW-U09-R04
% key: KEY-NEW-U09-R04
% mdparagraph:P294:begin
\textbf{R04. Complete controlled-reading translation.}
% mdparagraph:P294:end

% mdparagraph:P295:begin
The fisherman and his companion stand on the shore. The companion says, “Tell me what you have.” The fisherman says that he has food. He commands his companion to carry the fish. The companion carries the fish. The fisherman says, “I see you carrying the fish. Now I know that you are helping me.”
% mdparagraph:P295:end

% mdparagraph:P296:begin
\textbf{Target and interpretation:} Preserve all speakers and the progression from question to report, instruction, carrying and observation. \textbf{Hē/\allowbreak{}his} continues the fisherman; his companion carries one fish. \textbf{Birð} is simple present, without a completed-action claim. “Stand at the shore,” “tells his companion to carry” and “I see you carry” are defensible. Natural English need not visibly mark the subjunctive, but it must retain reported “says” versus asserted “know.”
% mdparagraph:P296:end

\end{minipage}\par\addvspace{4pt}

\reserve{7\baselineskip}
% mdparagraph:P297:begin
\section*{S. The Choice of Occupations}
% mdparagraph:P297:end

% mdparagraph:P298:begin
Use the eleven local helps supplied with the exact source. These answers interpret the printed passage; they do not turn every displayed source form into an independent production target.
% mdparagraph:P298:end

\par\begin{minipage}{\linewidth}\raggedright
% item: NEW-U09-S01
% key: KEY-NEW-U09-S01
% mdparagraph:P299:begin
\textbf{S01. Two comparison domains.}
% mdparagraph:P299:end

% mdparagraph:P300:begin
\textbf{Godes ðēowdōm}, service of God, receives first place overall. \textbf{Eorðtilð}, agriculture or farming, receives first place among worldly occupations.
% mdparagraph:P300:end

% mdparagraph:P301:begin
\textbf{Target and interpretation:} The third turn narrows the comparison to \textbf{woruldcræftas}. The second answer therefore does not withdraw the first. Making agriculture the first answer, or treating God's service as the answer within the later restricted worldly comparison, loses the exchange's structure.
% mdparagraph:P301:end

\end{minipage}\par\addvspace{4pt}

\par\begin{minipage}{\linewidth}\raggedright
% item: NEW-U09-S02
% key: KEY-NEW-U09-S02
% mdparagraph:P302:begin
\textbf{S02. The Gospel quotation.}
% mdparagraph:P302:end

% mdparagraph:P303:begin
The quotation begins \textbf{Fyrmest sēceað rīce Godes} and ends \textbf{tōgēiecte ēow}: “First seek God's kingdom and his righteousness, and all these things will be added to you.”
% mdparagraph:P303:end

% mdparagraph:P304:begin
It commands the hearers to seek the kingdom and righteousness first; it promises the addition of all these things. \textbf{Sēceað} is a plural imperative.
% mdparagraph:P304:end

% mdparagraph:P305:begin
\textbf{Target and interpretation:} The command addresses several hearers. \textbf{Ðās ðing eall} is the subject of indicative \textbf{bēoð}; dative \textbf{ēow} names recipients. The supplied passive promises a result, not another imperative ordering them to add things. Preserve the actual source \textbf{sēceað} beside the selected productive \textbf{sēcað}, rather than silently rewriting the quotation.
% mdparagraph:P305:end

\end{minipage}\par\addvspace{4pt}

\par\begin{minipage}{\linewidth}\raggedright
% item: NEW-U09-S03
% key: KEY-NEW-U09-S03
% mdparagraph:P306:begin
\textbf{S03. Two datives, different roles.}
% mdparagraph:P306:end

% mdparagraph:P307:begin
In \textbf{Ic sęcge ðē}, dative \textbf{ðē} is the speech recipient: the Counsellor addresses the Teacher. In \textbf{mē is geðūht}, dative \textbf{mē} is the experiencer, “to me” in “it seems to me”: the judgment is the Counsellor's.
% mdparagraph:P307:end

% mdparagraph:P308:begin
\textbf{Target and interpretation:} The first dative identifies who hears; the second identifies whose judgment follows. Swapping them gives the Teacher the Counsellor's judgment. The complete seems expression remains supplied; its past participle does not turn the present expression into “I thought.”
% mdparagraph:P308:end

\end{minipage}\par\addvspace{4pt}

\par\begin{minipage}{\linewidth}\raggedright
% item: NEW-U09-S04
% key: KEY-NEW-U09-S04
% mdparagraph:P309:begin
\textbf{S04. Why agriculture?}
% mdparagraph:P309:end

% mdparagraph:P310:begin
\textbf{For-ðām se ierðling ūs ealle fēt} — because the farmer feeds us all.
% mdparagraph:P310:end

% mdparagraph:P311:begin
\textbf{Se ierðling} is the nominative singular subject of \textbf{fēt}; \textbf{ūs ealle} is the accusative plural object, all of us.
% mdparagraph:P311:end

% mdparagraph:P312:begin
\textbf{Target and interpretation:} The farmer supplies food to the speakers and community; \textbf{ūs} includes the speaker but gives no exhaustive list of individuals. Plural \textbf{ealle} describes those fed, not the singular farmer. Exact \textbf{fēt} remains supplied from \textbf{fēdan}, not inferred productively by removing the prefix from \textbf{āfēdan}.
% mdparagraph:P312:end

\end{minipage}\par\addvspace{4pt}

\par\begin{minipage}{\linewidth}\raggedright
% item: NEW-U09-S05
% key: KEY-NEW-U09-S05
% mdparagraph:P313:begin
\textbf{S05. Complete source translation.}
% mdparagraph:P313:end

% mdparagraph:P314:begin
\textbf{Teacher:} “What do you say, wise man? Which occupation seems to you to be foremost among these?”
% mdparagraph:P314:end

% mdparagraph:P315:begin
\textbf{Counsellor:} “I tell you, God's service seems to me to hold first place among these occupations, as it is read in the Gospel: ‘First seek God's kingdom and his righteousness, and all these things will be added to you.’”
% mdparagraph:P315:end

% mdparagraph:P316:begin
\textbf{Teacher:} “And which seems to you to hold first place among worldly occupations?”
% mdparagraph:P316:end

% mdparagraph:P317:begin
\textbf{Counsellor:} “Agriculture, because the farmer feeds us all.”
% mdparagraph:P317:end

% mdparagraph:P318:begin
\textbf{Target and interpretation:} Retain all four turns, the complete quotation and the narrowed third-turn comparison. “Craft/\allowbreak{}trade,” “chief place/\allowbreak{}primacy,” “farming/\allowbreak{}tillage” and “plowman” are defensible equivalents. “It seems to me that God's service holds first place” is also natural English. God is the object of the service; the service, not God, receives first place.
% mdparagraph:P318:end

% mdparagraph:P319:begin
Translate \textbf{is geðūht} as a present judgment. Cook identifies \textbf{furðra} as nominative singular; the first question's \textbf{hwilc} is also nominative. Do not force the whole source frame into the command/\allowbreak{}perception construction with an accusative actor.
% mdparagraph:P319:end

% mdparagraph:P320:begin
Cook's note 12 marks three \textbf{is} insertions as conjectural, absent from his manuscripts: those in the two Teacher questions and the \textbf{swā-swā} clause. The distinct \textbf{mē is geðūht} is not a fourth application of that note. This translation follows his printed text while retaining that qualification. Note 20 records manuscript \textbf{cræftas woruld} beside printed \textbf{woruldcræftas}. The printed compound, participle \textbf{tōgēiecte} and command \textbf{sēceað} remain as given, with local helps. [4]
% mdparagraph:P320:end

\end{minipage}\par\addvspace{4pt}

\reserve{7\baselineskip}
% mdparagraph:P321:begin
\section*{X. Connected application and repair}
% mdparagraph:P321:end

\par\begin{minipage}{\linewidth}\raggedright
% item: NEW-U09-X01
% key: KEY-NEW-U09-X01
% mdparagraph:P322:begin
\textbf{X01. A connected account.}
% mdparagraph:P322:end

% mdparagraph:P323:begin
\textbf{Se fiscere hēt ðā męnn ðone fisc beran. Ic geseah hīe hine beran. Ic wāt, ðæt hīe him hulpon.}
% mdparagraph:P323:end

% mdparagraph:P324:begin
The fisherman commanded the men to carry the fish. I saw them carry it. I know that they helped him.
% mdparagraph:P324:end

% mdparagraph:P325:begin
\textbf{Target and interpretation:} \textbf{Ðā męnn/\allowbreak{}hīe} are the commanded and observed plural actors; \textbf{ðone fisc/\allowbreak{}hine} is the singular fish. Masculine \textbf{hine} follows \textbf{fisc} although English says “it.” In the final clause \textbf{hīe} is nominative subject of preterite indicative \textbf{hulpon}, and \textbf{him} the dative person helped, naturally the fisherman. Keep present \textbf{wāt} with factual past help; subjunctive \textbf{hulpen} misses the target. Plain \textbf{beran} gains neither plural nor past endings.
% mdparagraph:P325:end

\end{minipage}\par\addvspace{4pt}

\par\begin{minipage}{\linewidth}\raggedright
% item: NEW-U09-X02
% key: KEY-NEW-U09-X02
% mdparagraph:P326:begin
\textbf{X02. Repairing deliberate errors.}
% mdparagraph:P326:end

% mdparagraph:P327:begin
\textbf{(a) Ic gesēo hine ðone fisc beran.} — I see him carrying the fish. Replace nominative \textbf{hē} with accusative actor \textbf{hine}.
% mdparagraph:P327:end

% mdparagraph:P328:begin
\textbf{(b) God geworhte ðone mann gōdne.} — God made the man good. Replace \textbf{gōd} with strong masculine accusative singular \textbf{gōdne}, agreeing with the object.
% mdparagraph:P328:end

% mdparagraph:P329:begin
\textbf{Target and interpretation:} Preserve the selected complements while repairing case. Being the carrier does not make the expressed actor nominative after \textbf{gesēo}. The predicate adjective describes the object, so its ending follows that accusative object, not nominative subject \textbf{God}. The erroneous inputs were deliberately composed for diagnosis, not source readings to preserve.
% mdparagraph:P329:end

\end{minipage}\par\addvspace{4pt}

\clearpage
% mdparagraph:P330:begin
\section*{Source and editorial notes}
% mdparagraph:P330:end

\begingroup\fontsize{10}{12.8}\selectfont
% mdparagraph:P331:begin
[1] Cook, \emph{A First Book in Old English}, second edition (1895), §§169, 173 and 194–199, supplies the central complement constructions; Fulk, \emph{An Introductory Grammar of Old English} (2020), §§44–45 and 58, supports the stated mood and infinitive distinctions. Dependent-question frames were also checked against Bosworth–Toller, Main volume p. 855b, \textbf{secgan}, and Supplement p. 575a–b, \textbf{hwā}. These bounded assignments remain separate from later tense-sequence, general impersonal and additional passive systems.
% mdparagraph:P331:end

% mdparagraph:P332:begin
[2] The selected new noun paradigms follow Cook, \emph{First Book}, printed pp. 28–29, 31 and 33–34; the verbal systems follow the audited paradigms and lexical entries. Bosworth–Toller, Main volume p. 137, corroborates the selected \textbf{bycgan} forms. Documented alternatives remain valid; a requested selected table is a narrower exercise target.
% mdparagraph:P332:end

% mdparagraph:P333:begin
[3] Cook, \emph{Exercises in Old English} (1895), Group II, Exercises 9–12, printed pp. 4–6 /\allowbreak{} PDF pp. 12–14: all 64 prompts. Parentheses, brackets and the final comma in II.9.11 belong to those prompts. The Old English model answers are new teaching material. Absolute \textbf{slēan} is supported by Bosworth–Toller, Main volume p. 883a, B.I; Cook, \emph{First Book}, p. 132 also supports a hammer-striking use of \textbf{bēatan}. The numeral help in II.12.13 is an explicit adaptation of Cook §§78–79, not a claimed literal source sentence.
% mdparagraph:P333:end

% mdparagraph:P334:begin
[4] The exact four-turn source is Cook, \emph{First Book}, printed p. 131 /\allowbreak{} PDF p. 147, prose lines counted editorially as 7–15, with printed 10 and 15 as anchors. The complete translation includes its Gospel quotation, but not the preceding merchandise exchange or following smith speech. All thirteen original notes remain in the learner chapter. Notes 12, 14 and 20 govern the conjectural copulas, nominative \textbf{furðra} and manuscript word order discussed in S05. Unfamiliar source constructions retain the eleven local helps. PDF pages here are one-based.
% mdparagraph:P334:end

\endgroup
\end{document}
